% =====================================================================
% §Limitations and Honest Disclosures  [F12]
% Fragment \input inside \section{Limitations and Honest Disclosures}
% All numbers from _DRAFTING_BRIEF.md (authoritative).
% =====================================================================

The framework is validated across the component benchmarks of
Sec.~\ref{sec:benchmarks} and the externally selected 25-flight corpus of
Sec.~\ref{sec:corpus}, but several limitations bound the scope of admissible
inference. Each is disclosed below with its quantitative magnitude, its root
cause where identified, and either an in-paper mitigation or a documented
future fix. We state these plainly so that the headline claim---supersonic
parity with RASAero~II to $M = 4.33$, with hypersonic reach demonstrated only
as an exploratory method capability---is not over-read.

\subsection{Transonic regime weakness}
\label{sec:lim-transonic}

The largest disclosed model weakness is in the transonic band. Across the
$n = 7$ headline flights with peak $M = 0.8$--$1.3$, OpenRocket-Plus carries a
mean signed apogee error of $-3.66\%$ ($\sigma = 5.36\%$, RMSE $6.17\%$,
MAE $5.84\%$), the only regime in which RASAero~II~\cite{rogers2015} is the more
accurate predictor. The likely cause is that the supersonic-tuned blending of
the wave-drag and base-drag stacks pulls total drag slightly low on the subsonic
side of $M = 1$, over-predicting apogee. The transonic pitch-damping
augmentation provides an independent corroboration of the same calibration gap:
the Gaussian $C_{mq}$ multiplier overshoots the transonic damping peak by $+110$
to $+160\%$ at $M = 1.08$--$1.11$ relative to the Sznajder Fluent
reference~\cite{sznajder2025}. Two independent observations---flight-corpus
apogee bias and CFD-benchmarked $C_{mq}$---therefore converge on the transonic
blending region as the principal open calibration gap. The proposed fix is to
redistribute transonic component wave drag through a Whitcomb equivalent body of
revolution rather than summing the components independently.

\subsection{Partly in-sample headline (two B-level base-drag constants)}
\label{sec:lim-insample}

The headline corpus is \emph{partly} in-sample. Two base-drag scale constants
are B-level and corpus-frozen: a thick-boundary-layer factor
$\mathtt{THICK\_BL\_K} = 2.2$ (anchored on a single flight, Raven) and a
slender-body factor $\mathtt{SLENDER\_BODY\_K} = 0.0025$ (anchored on
Raven, Rabia, and Kinsel, with Torrent used in the source diagnostic). Because
these constants were tuned against flights that also appear in the headline
statistics, the headline cannot be read as fully blind. The mitigation is a
\emph{decontaminated} prospective holdout in which every flight any constant
touched (Raven, Rabia, Rabia Short Fin Can, Kinsel, and Torrent) is placed in
the development partition. On that split the development partition ($n = 13$)
gives a mean error of $+0.22\%$ and MAE $5.47\%$, while the genuinely blind
holdout ($n = 12$) gives a mean of $-1.03\%$ and MAE $3.95\%$. The holdout is
\emph{more} accurate than the development set, which is the primary evidence
that the two constants generalize rather than overfit. We note that an earlier
split inadvertently placed a constant-touched flight (Rabia) in the holdout; the
decontaminated split reported here corrects that contamination.

\subsection{Sparse high-Mach sampling in the headline corpus}
\label{sec:lim-sparse}

The headline corpus peaks at $M = 4.33$; its hypersonic regime ($M > 5$) is
\emph{empty}. Of the 25 flights, 9 are subsonic ($M < 0.8$), 7 transonic, 5 low
supersonic ($M = 1.3$--$3.0$), and only 4 high supersonic ($M = 3.0$--$5.0$).
Above $M = 3$ the per-regime statistics are descriptive only ($n \le 5$); no
inferential claim is attached to them. Consequently the hypersonic capability is
reported in Sec.~\ref{sec:exploratory} strictly as an \emph{exploratory}
result---the method \emph{reaches} $M = 7$ within $\pm 7\%$ on
well-characterized vehicles (Black~Brant~VB at $M = 7.22$, $-6.97\%$; two
Nike-Deacon flights at $M \approx 5$, $-1.06\%$ and $-0.89\%$)---but on poorly
documented historical flights the motor- and geometry-reconstruction uncertainty
dominates, and the method is \emph{not} validated in the hypersonic regime. We
deliberately do not present the high-Mach set as a within-tolerance headline,
which would constitute outcome-based selection.

\subsection{Two-stage closure (MESOS 293K) regression}
\label{sec:lim-mesos}

The single two-stage closure in the headline corpus, MESOS~293K, currently
predicts a signed apogee error of $-6.96\%$. An archived value of $-0.6\%$ is
\emph{not} reproducible from the present code and is disclosed here as a
regression rather than carried forward. The current $-6.96\%$ result remains
within the $\pm 10\%$ admission criterion and so does not change the
$25/25$ headline, but the discrepancy is acknowledged and the cause---most
plausibly staging-event mass and drag bookkeeping across the interstage---is
flagged for diagnosis. We report the reproducible $-6.96\%$ value throughout.

\subsection{No author-produced computational fluid dynamics}
\label{sec:lim-cfd}

The validation in Sec.~\ref{sec:cfd} cites four published-CFD comparators rather
than runs produced by the present author; this is a deliberate scoping decision
for a solo-author, self-funded open-source contribution. The mitigation is
breadth: the comparators span two geometries (Basic~Finner and AGARD-B), two
coefficient families (static drag and dynamic pitch damping), the transonic
through hypersonic-limit Mach bands, and four independent author
groups---Bunescu URANS~\cite{bunescu2025}, the Sahu and Bhagwandin TLNS and
follow-up work~\cite{sahu1983,bhagwandin2013}, Vidanovi\'c SST
$k$--$\omega$~\cite{vidanovic2014}, and Sznajder Fluent~\cite{sznajder2025}.
Future work is to ship a closed-loop comparator pipeline---an author-built
AGARD-B model driving the Vidanovi\'c reference and an ogive-cylinder-boattail
model driving the Sahu base-drag reference---to remove the ``reference dataset
only'' caveat.

\subsection{B-level transonic pitch-damping multiplier}
\label{sec:lim-cmq}

The transonic pitch-damping augmentation is held at B-level confidence. The
strip-theory damping model applies a single transonic multiplier of
approximately $3\times$, calibrated against AEDC pitch-damping
data~\cite{aedc7658} and re-checked against the Sznajder Fluent
reference~\cite{sznajder2025}, against which it overshoots the transonic damping
peak by $+110$ to $+160\%$ at $M = 1.08$--$1.11$ (Sec.~\ref{sec:lim-transonic}).
The root cause is the absence of a flight-scale anchor: the multiplier rests on
a single wind-tunnel-scale source with no independent full-scale validation.
This affects predicted dynamic stability and coning behavior but not the apogee
statistics, which are insensitive to $C_{mq}$. Lifting this term to A-level
requires an independent flight-scale pitch-damping dataset, which is identified
as future work.

\subsection{Heterogeneous ground-truth sources}
\label{sec:lim-truth}

The corpus ground truth is heterogeneous, drawn from GPS- and barometric-altimeter
apogee records, optical and radar tracking, and published flight reports of
differing vintage and instrumentation. Reported apogee precisions across these
sources imply an irreducible measurement floor on the order of $1\%$ that is
\emph{not} subtracted from the reported errors. The headline mean signed error
of $-0.38\%$ and RMSE of $5.34\%$ therefore include this measurement
contribution; the true model-only error is correspondingly smaller, but we make
no attempt to deconvolve it and report the raw, conservative figures throughout.

\subsection{Aeroelastic model implemented but disabled}
\label{sec:lim-aeroelastic}

A fin aeroelastic-effectiveness model is present in the codebase but is gated by
a dynamic-pressure threshold of $Q_{\mathrm{THRESHOLD}} = 1 \times 10^{12}$~Pa
and is therefore effectively disabled in every simulation reported here. No
aeroelastic claims---flutter, divergence, or fin-tip-twist-driven $C_{N\alpha}$
reduction---are made in this paper. The threshold is held at $10^{12}$~Pa
pending an author flutter-and-divergence validation campaign against published
cantilever-fin data, after which the gate can be lowered to a physically
meaningful value.

\subsection{Sensitivity sweep run on the exploratory high-Mach vehicles}
\label{sec:lim-sensitivity}

The four-flight operational sensitivity sweep of Sec.~\ref{sec:corpus} was run on
the high-Mach \emph{exploratory} vehicles (HEROS~3, the blunt Arcas, a Nike-Apache,
and Black~Brant~VB) rather than on the gentler headline flights. These are the
fastest, highest-Mach trajectories in the dataset and therefore bound the
numerical-convergence behavior of the milder headline cases as a worst-case
envelope: drag-coefficient scale dominates the response (mean local sensitivity
$|s| \approx 4\%$ apogee per $10\%$ $C_d$, i.e.\ a sensitivity coefficient near
$0.4$), while the integration time step contributes $|s| \approx 0.98\%$ over the
range $0.025$--$0.10$~s, confirming numerical convergence per the AIAA accuracy
policy~\cite{aiaa_numerical_policy}. We state explicitly that the sweep vehicles
are exploratory: the in-sample defense of the headline rests on the
decontaminated holdout of Sec.~\ref{sec:lim-insample}, not on this sweep.

\begin{figure}[t]
  \centering
  \includegraphics[width=\columnwidth]{sensitivity_tornado_composite}
  \caption{Sensitivity tornado composite for the four high-Mach exploratory
  vehicles (HEROS~3, blunt Arcas, Nike-Apache, Black~Brant~VB) under
  $\pm 10\%$ perturbations of drag-coefficient scale, integration time step,
  launch altitude, and launch-rod angle. These worst-case trajectories bound the
  numerical-convergence behavior of the gentler headline flights; drag-coefficient
  scale dominates the response.}
  \label{fig:sensitivity-tornado}
\end{figure}
